\section{Discussion}
\label{sec:discussion}

\paragraph{The INT8 latency paradox, quantified.} INT8 quantization divides the weight memory footprint by
4 but does \emph{not} speed up inference on the Cortex-M4 --- it slows it down. Cycle-level instrumentation of
the kernel (Section~\ref{sec:system}) says where: across both datasets, the integer MAC costs on average
\num{6785}~cycles and requantization \num{3777.25}, while dequantization weighs only \num{200}. The INT8
total settles at \SI{74}{\micro\second} against \SIrange{48}{50}{\micro\second} in FP32. Two causes therefore
compound: the integer MAC has nothing to gain against a hardware FPU already efficient in single precision,
and requantization --- an item that does not exist in the floating-point path --- consumes a third of the
budget on its own. The direct lead remains a port to integer \gls{SIMD} kernels, such as those of
\gls{CMSISNN}~\cite{Lai2018CMSISNN}, which would attack the
first item; the second rather calls for a power-of-two scale chain.
This balance concerns inference. Under continual learning the overhead is heavier still, by a factor
of~\num{2.4} (Section~\ref{sec:recovery-board}), since rebuilding the integer view follows every gradient
step. That is the regime on-device learning claims, and therefore where the paradox weighs most.

\paragraph{An energy paradox compounds the latency paradox.} One could hope that INT8, failing to be faster,
would at least be leaner. Measurement says otherwise: \SI{67.47}{\micro\joule} per inference in INT8 against
\SI{67.65}{\micro\joule} in FP32, a gap of \SI{0.18}{\micro\joule} for a combined uncertainty of
\SI{2.34}{\micro\joule}. There is thus no measurable effect. The practical consequence is clear: on this
target, \textbf{INT8 is justified by memory, neither by speed nor by energy}. The energy lever lies
elsewhere --- in clock frequency, where the ample latency headroom is exchanged for battery life.

\paragraph{"Measured" vs "emulated" honesty.} We rigorously separate two kinds of results. FP32 parities, the
legacy INT8 collapse, recovery through the v2 kernel, latencies, RAM and energy are \emph{measured on real
hardware} (Sprints~36, 40, 46 to 50 and 53). The causal diagnosis, the ablation ladder and the depth sweep
are \emph{bit-exactly emulated on PC} (Sprints~39 and~47): the emulator reproduces the C kernel arithmetic
exactly, which makes it a reliable reference, but it is not a board measurement. Since the v2 campaign,
recovery no longer belongs to the second category: it is established on the four per-channel cells, with a
frozen parity of \num{1.000} against the emulator. What remains unmeasured --- the \texttt{q15} format
on-board, energy per continual update, the per-phase profile --- is reported as "to be measured" rather than
estimated.

\paragraph{Limitation: forgetting is not stressed on these two datasets.} Our continual scenarios there are
short (three tasks) and average forgetting stays below \num{0.01}: these datasets do not put EWC
regularization to the test. Catastrophic forgetting was measured elsewhere in the project, on CWRU, where it
reaches \num{0.8475} for EWC against \num{0.8578} under naive learning. As that dataset does not overlap
ours, we refrain from mixing the two findings: this article's contribution bears on \emph{porting} and
\emph{quantization}, not on the effectiveness of regularization against forgetting.

\paragraph{Scope.} The central result --- a naive PTQ can destroy a binary decision head, and a
per-tensor/per-channel calibration (or a Q15 format) restores it --- extends beyond the EWC case: it concerns
any lightweight head quantized downstream of a feature extractor, a frequent regime in TinyML.
